\documentclass[a4paper]{article}
\usepackage{amsfonts}
\usepackage{amsmath}
\usepackage{amssymb}
\usepackage{graphicx}     

\begin{document}

\noindent \large Auteur: Alhassane Balde\\
\noindent \large Studierichting: Informatica\\
\noindent \large Oefeningenreeks 2E nr. 9.5\\

\medskip

\normalsize


$ \left( (3^x)^x \right)^{x+1} \cdot 3^{-x} = 3 $ \\  

$ \Rightarrow 3^{x \cdot x \cdot (x+1)} \cdot 3^{-x} = 3 $ \\  

$ \Rightarrow 3^{x^2(x+1) - x} = 3^1 $ \\ 

$ \Rightarrow x^2(x+1) - x = 1 $ \\  

$ \Rightarrow x^3 + x^2 - x - 1 = 0 $ \\

We lossen de vergelijking $ x^3 + x^2 - x - 1 = 0 $ op. \\  

We zoeken een gehele oplossing door de delers van $ -1 $ te testen: $ x \in \{\pm 1\} $. \\  

Voor $ x = 1 $: \\  

$ 1^3 + 1^2 - 1 - 1 = 1 + 1 - 1 - 1 = 0 $ \\  

Dus $ x = 1 $ is een oplossing. \\  

We delen $ x^3 + x^2 - x - 1 $ door $ (x - 1) $ met polynomiale deling: \\  

$ (x - 1)(x^2 + 2x + 1) = 0 $ \\  

Nu lossen we $ x^2 + 2x + 1 = 0 $ op: \\  

$ (x + 1)^2 = 0 $ \\  

$ x + 1 = 0 \Rightarrow x = -1 $ \\  

De oplossingen van de vergelijking zijn: \\  

$ x = 1 \quad \text{of} \quad x = -1 $  


\begin{thebibliography}{999}
    %\bibitem{a} Referentie
\end{thebibliography}
\end{document}
